\chapter{Related Work}
\label{ch:related-work}

\section{Formal Methods and Automaton-Based ICS Detection}
\label{sec:related-formal}

Model checking and automaton-based verification have a long history in protocol
analysis, with the foundational reference text by Clarke and colleagues setting the
agenda for several decades of work \autocite{ClarkeModelChecking2001}. In the ICS
domain, the NES@FIT group at the Faculty of Information Technology, Brno University
of Technology has produced a series of automaton-based detectors for IEC-104
communication \autocite{HolikMatousekRysavyVojnar2023}. The output of these
detectors is a structured verdict with class, affected devices, and timing
information, and it forms the formally verified bound on which the work of this
thesis is grounded.

A complementary line of work in the first supervisor's previous publications has
addressed formal reachability analysis in networks with dynamic behaviour
\autocite{DeSilva2014FormalReachability}. The methodology there showed that the
output of a formal analysis can be structured for downstream consumption rather than
treated solely as an internal verification artefact; the same intuition motivates
this thesis.

[Author's note: extend this section with the M2 reading notes; emphasise the gap
between the formal-output side and the AI-explanation side that this thesis closes.]

\section{LLM Faithfulness and Hallucination Research}
\label{sec:related-faithfulness}

The faithfulness of large-language-model outputs has been studied extensively in the
context of summarisation \autocite{Maynez2020Faithfulness} and in long-form text
generation \autocite{Min2023FActScore}. Honovich and colleagues established a common
evaluation methodology for factual consistency by re-aligning several existing
benchmarks \autocite{Honovich2022True}. These contributions share a common
methodology: the output of the model is compared against a reference (a source
document or a knowledge base) and a score is produced.

The framework presented in this thesis differs from these in two ways. First, the
reference is not a source document but a formally verified detector output, and the
comparison is therefore against a machine-readable bound rather than against natural
language. Second, the comparison is performed at runtime so that violations can be
acted on rather than only reported.

[Author's note: extend this section with the M2 reading notes; show how FActScore-style
claim decomposition informed the per-claim three-state classifier described in
Chapter~\ref{ch:framework}.]

\section{Runtime Verification Theory}
\label{sec:related-rv}

Runtime verification is the formal area concerned with checking whether an execution
of a system satisfies a specified property at run time
\autocite{LeuckerSchallhart2009RuntimeVerification}. The classical setting checks
program executions against linear-temporal-logic or related specifications. This
thesis applies the runtime-verification viewpoint to AI explanations: the property to
be verified is the faithfulness of the explanation against the formal detector
output, and the monitor specification is the bound itself.

[Author's note: discuss the difference between safety properties and liveness
properties in the context of explanations; the bounds described in
Chapter~\ref{ch:framework} are safety properties.]

\section{The European Union AI Act and Safety-Critical AI}
\label{sec:related-eu-ai-act}

The European Union AI Act introduces transparency obligations (Article~13) and
quality-management obligations (Article~17) for AI systems deployed in critical
infrastructure (Annex~III). These provisions are particularly relevant for AI systems
that translate the output of safety-critical detection mechanisms into natural
language for operators, since both the form and the content of the explanation become
regulated outputs.

[Author's note: the operational deployment guidelines presented in
Chapter~\ref{ch:discussion} are mapped to these provisions; cite the AI Act consolidated
text and the relevant Commission guidance.]

\section{Position of This Thesis}
\label{sec:position}

The four lines of work surveyed above each contribute one component of the framework
introduced here. Formal methods and automaton-based detection contribute the verified
bound layer. Faithfulness research contributes the methodology for decomposing an
explanation into atomic claims. Runtime verification contributes the theoretical
foundation for evaluating a property of an execution at run time. The EU AI Act
contributes the governance context that motivates an operational guideline.

The contribution of this thesis is the integration of these four threads into a
single runtime verification framework for AI reasoning that is grounded in formally
verified detector outputs and that is empirically validated in the ICS domain.
